\section{Datasets y fuentes}
\label{sec:datasets}

Se usan cinco fuentes oficiales y públicas (Tabla~\ref{tab:datasets}): cuatro de la Plataforma Nacional de
Datos Abiertos \cite{datosabiertos} y una de la API de interoperabilidad del SENAMHI. La llave común es el
\code{ubigeo}, el código de 6 dígitos de cada distrito (2 del departamento, 2 de la provincia y 2 del distrito).

\begin{table}[htbp]
\centering
\caption{Datasets del proyecto. \emph{Filas}: tamaño de cada tabla después de la preparación con Spark.}
\label{tab:datasets}
\small
\begin{tabularx}{\textwidth}{@{}>{\raggedright\arraybackslash}p{4.1cm} >{\raggedright\arraybackslash}p{2.2cm} X r@{}}
\toprule
Dataset (publicador) & Cobertura & Contenido y uso & Filas \\
\midrule
Vigilancia de dengue (CDC Perú – MINSA) \cite{ds-dengue} & 2000–2025 &
  Un caso notificado por fila: distrito, semana epidemiológica, edad, sexo y diagnóstico (CIE-10).
  Tabla principal. & 1,065,274 \\
Población identificada con DNI (RENIEC) \cite{ds-reniec} & dic. 2022–2025 &
  Población por distrito, sexo y edad. Denominador de las tasas. & 75,636 \\
Distritos en emergencia hídrica (SUNASS) \cite{ds-sunass} & 2017, 2023, 2024 &
  Un distrito por decreto: tipo de emergencia (lluvias o déficit hídrico) y horas de agua al día. & 9,538 \\
Establecimientos de salud, RENIPRESS (SUSALUD) \cite{ds-renipress} & corte al 30-09-2026 &
  Establecimientos activos por distrito y nivel de atención. & 26,972 \\
Estaciones meteorológicas (SENAMHI) \cite{senamhi-api} & 2022–2025 &
  Lluvia, temperatura y humedad diarias de 9 estaciones de la costa norte y la selva. & 12,827 \\
\bottomrule
\end{tabularx}
\end{table}

La población por distrito solo existe desde 2022, así que todas las tasas usan 2022–2025. Es la población con
DNI según el domicilio declarado, no la población censal. Los datos del SENAMHI no tienen control de calidad y
algunas estaciones tienen días sin registro.

Para que el análisis sea reproducible, \code{data/catalog.yaml} registra la URL y el hash SHA-256 de cada
archivo, y un script los descarga y verifica. Los datos no se versionan: solo el CSV de dengue pesa 118~MB.
